% Part: normal-modal-logic
% Chapter: frame-correspondence
% Section: first-order-definability

\documentclass[../../../include/open-logic-section]{subfiles}

\begin{document}

\olfileid{nml}{frd}{fol}

\olsection{మొదటిస్థాయి నిర్వచనీయత}

చట్రాల ప్రాప్యత సంబంధాలకు ఉన్న అనేక ధర్మాలను మోడల్
\tetoken{సూత్రాలతో}{!!{formula}s} నిర్వచించవచ్చని చూశాం. ఉదాహరణకు,
చట్రాల సౌష్ఠవాన్ని \Ax{B} అనే \tetoken{సూత్రం}{!!{formula}},
$p \lif \Box\Diamond p$, నిర్వచిస్తుంది. ఇప్పటివరకు ఎదురైన
షరతులన్నింటినీ ఒక్క ద్విస్థాన \tetoken{విధేయ సంకేతం}{!!{predicate}}
ఉన్న \tetoken{భాషలోని}{!!a{language}} మొదటిస్థాయి
\tetoken{సూత్రాలతో}{!!{formula}s} వ్యక్తం చేయవచ్చు. ఉదాహరణకు,
$\lforall[x][\lforall[y][(\Atom{Q}{x,y} \lif \Atom{Q}{y, x})]]$
సౌష్ఠవాన్ని నిర్వచిస్తుంది: $\Domain{M} = W$,
$\Assign{Q}{M} = R$ అయ్యే మొదటిస్థాయి
\tetoken{నిర్మాణం}{!!{structure}}~$\Struct{M}$ ఆ సూత్రాన్ని
సంతృప్తిపరచేది $R$ సౌష్ఠవమైతేనే. ఇది కింది నిర్వచనాన్ని
సూచిస్తుంది:

\begin{defn}
  చట్రాల వర్గం~$\mClass{F}$ను \emph{మొదటిస్థాయి
    నిర్వచనీయమైనది} అంటాం, ఒకవేళ ఒక్క ద్విస్థాన
  \tetoken{విధేయ సంకేతం}{!!{predicate}}~$Q$ ఉన్న
  మొదటిస్థాయి భాషలో ఒక \tetoken{వాక్యం}{!!a{sentence}}~$!A$
  ఉండి, ప్రతి చట్రం విషయంలో
  $\mModel{F} = \tuple{W, R} \in \mClass{F}$ అనే షరతు,
  $\Domain{M} = W$, $\Assign{Q}{M} = R$ గల
  మొదటిస్థాయి \tetoken{నిర్మాణం}{!!{structure}}~$\Struct{M}$లో
  $\Sat{M}{!A}$ అనే షరతుకు సమానమైతే.
\end{defn}

ఇప్పటివరకు పరిశీలించిన ధర్మాలూ, వాటిని నిర్వచించే మోడల్
\tetoken{సూత్రాలూ}{!!{formula}s} ప్రత్యేకమైనవని తేలుతుంది.
ప్రతి మోడల్ \tetoken{సూత్రం}{!!{formula}} మొదటిస్థాయి
నిర్వచనీయమైన చట్రాల వర్గాన్ని నిర్వచించదు; మొదటిస్థాయి
నిర్వచనీయమైన ప్రతి చట్రాల వర్గాన్నీ మోడల్ సూత్రంతో
నిర్వచించలేం.

మొదటి వాదనకు లొబ్ సూత్రం ప్రతిదృష్టాంతం:
\begin{equation}
\Box(\Box p \lif p) \lif \Box p. \tag{\Ax{W}}
\end{equation}
\Ax{W} సంక్రామకమూ విలోమంగా సుస్థాపితమూ అయిన చట్రాల
వర్గాన్ని నిర్వచిస్తుంది. $Rw_2w_1$, $Rw_3w_2$, \dots{} అయ్యే
$w_1$, $w_2$, \dots{} అనే అనంత శ్రేణి లేకపోతే ఆ సంబంధాన్ని
\emph{సుస్థాపితమైనది} అంటాం. ఉదాహరణకు, $\Nat$పై $<$
సుస్థాపితమైనది; $\Int$పై $<$ అలా కాదు. ఒక సంబంధం
\emph{విలోమంగా సుస్థాపితమైనది} అయ్యేది దాని విలోమం
సుస్థాపితమైనప్పుడే. అందువల్ల $Rw_1w_2$, $Rw_2w_3$, \dots{}
అయ్యే $w_1$, $w_2$, \dots{} అనే అనంత శ్రేణి లేని సంబంధాలే
విలోమంగా సుస్థాపితమైనవి.

అయితే సంక్రామక, విలోమంగా సుస్థాపిత సంబంధాలను
నిర్వచించే మొదటిస్థాయి \tetoken{సూత్రం}{!!{formula}} లేదు.
అలా ఉందనుకొని, $\Sat{M}{!F}$ అయినప్పుడే, అప్పుడే
$R = \Assign{Q}{M}$ సంక్రామకమూ విలోమంగా సుస్థాపితమూ
అనుకుందాం. $!A_n$ అనేది కింది \tetoken{సూత్రం}{!!{formula}}
అనుకుందాం:
\[
(\Atom{Q}{a_1,a_2} \land \dots \land
    \Atom{Q}{a_{n-1},a_{n}})
\]
\sourcecorrection{OLTENMLFRDFOL-001}{మూలంలోని సూచికల
  అంచు సందర్భం స్పష్టంగా చెప్పలేదు: $!A_1$ను సత్యమైన శూన్య
  సంయోగంగా చదువుతాం; ఎంచుకున్న పరిమిత ఉపసమితిలో $!A_n$
  ఏదీ లేకపోతే ఒక మూలక నమూనా సరిపోతుంది. $n \ge 2$కు
  ముద్రించిన అనుసంధాన సూత్రాన్ని యథాతథంగా ఉంచాం.}
ఇప్పుడు ఈ \tetoken{సూత్రాల}{!!{formula}s} సమితిని పరిశీలిద్దాం:
\[
\Gamma = \{!F, !A_1, !A_2, \dots\}.
\]
$\Gamma$లోని ప్రతి పరిమిత ఉపసమితి సంతృప్తిపరచదగినది:
ఆ ఉపసమితిలోని $!A_k$ల సూచికలలో గరిష్ఠమైనది $k$
అనుకుందాం; $\Domain{M_k} = \{1, \dots, k\}$,
$\Assign{a_i}{M_k} = i$, $\Assign{Q}{M_k} = <$ అని
నిర్ణయిద్దాం. $\{1, \dots, k\}$పై $<$ సంక్రామకమూ
విలోమంగా సుస్థాపితమూ కాబట్టి $\Sat{M_k}{!F}$.
నిర్మాణం వల్ల అన్ని $i \le k$కు $\Sat{M_k}{!A_i}$.
మొదటిస్థాయి విధేయ తర్కపు సంహతత్వ సిద్ధాంతం ప్రకారం
$\Gamma$ను సంతృప్తిపరిచే ఏదో ఒక
\tetoken{నిర్మాణం}{!!{structure}}~$\Struct{M}$ ఉంది.
ఊహ ప్రకారం $\Sat{M}{!F}$ కనుక $\Assign{Q}{M}$
విలోమంగా సుస్థాపితమైనది. కానీ
$\Assign{a_1}{M}$, $\Assign{a_2}{M}$, \dots{}
విలోమ సుస్థాపితత్వం నిషేధించే తరహా అనంత శ్రేణిని
ఏర్పరుస్తాయని స్పష్టమే.

రెండవ వాదనకు ప్రతిదృష్టాంతం సార్వత్రికత ధర్మం:
ప్రతి $u$, $v$కు $Ruv$. సార్వత్రిక చట్రాలు
$\lforall[x][\lforall[y][\Atom{Q}{x,y}]]$ అనే
మొదటిస్థాయి సూత్రం ద్వారా నిర్వచనీయమైనవి.
అయితే అన్ని సార్వత్రిక చట్రాల్లో, వాటిలో మాత్రమే,
చెల్లుబాటు అయ్యే మోడల్ సూత్రం ఏదీ లేదు. స్వతంత్రంగా
ఆసక్తికరమైన ఒక ఫలితం దీనికి కారణం: సార్వత్రిక చట్రాలన్నింటిలో
చెల్లుబాటు అయ్యే సూత్రాల సమితి, స్వావర్తన, సౌష్ఠవ,
సంక్రామక చట్రాలన్నింటిలో చెల్లుబాటు అయ్యే సూత్రాల
సమితితో ఒకటే. సార్వత్రికం కాని
స్వావర్తన, సౌష్ఠవ, సంక్రామక చట్రాలు ఉన్నాయి; కనుక
అన్ని సార్వత్రిక చట్రాల్లో చెల్లుబాటు అయ్యే ప్రతి
\tetoken{సూత్రం}{!!{formula}} కొన్ని సార్వత్రికం కాని
చట్రాల్లోనూ చెల్లుబాటు అవుతుంది.

\end{document}
